\section{Every jet of the tensor residue multiplier and its surviving perimeter}
\label{sec:ftr}

This calculation evaluates the multiplier in TWB42--42a. It retains the
original complete quartet, its common multiplicity $m$, the order
$k=4\ell+1\geq9$, and every Taylor coefficient of $v_h=g/h$ used by the
original packet. In particular it does not replace the tensor residue
form by the cyclic residue form when their comparison is singular.

Write the four original roots and their full local factors as
\[
 z_{\epsilon\eta}=\tfrac12+(2\epsilon-1)\delta
                 +\mathrm i(2\eta-1)\gamma,\quad
 \epsilon,\eta\in\{0,1\},\qquad
 h_z(s)=\frac{h(s)}{(s-z)^m},\quad v_h(s)=\frac{g(s)}{h(s)}.
 \tag{FTR.1}
\]
Here $\delta,\gamma>0$, each displayed zero has its complete order $m$,
and $h$ contains precisely those four factors. Thus $h_z(z)$ and
$v_h(z)$ are nonzero. Set
\[
 r=k(m-1),\quad e=r+1,\quad
 \lambda_{ab}=\frac k2+(2a-k)\delta+\mathrm i(2b-k)\gamma,
 \quad
 \chi(S)=\prod_{a,b=0}^k(S-\lambda_{ab})^e,
 \quad \chi_{ab}(S)=\frac{\chi(S)}{(S-\lambda_{ab})^e}.
 \tag{FTR.2}
\]
These are the original sum centers and full annihilator; $q=e(k+1)^2$.
Let $\mathcal E=\mathbb C[S]/\chi$, with action $A$ given by
multiplication by $S$ and reflection
$f^{\sharp_k}(S)=\overline{f(k-\overline S)}$.
Use precisely the TWB cyclic residue
$R_\chi(f,g)=L_\chi(f^{\sharp_k}g)$, where
$L_\chi(P)=\sum_\lambda\operatorname{Res}_{S=\lambda}P(S)dS/\chi(S)$.
The original tensor pullback has the unique expression
\[
 B(f,g)=R_\chi(f,bg),\qquad b\in\mathcal E.
 \tag{FTR.3}
\]
Uniqueness and existence were proved in TWB42 by the actual cyclic
commutant. The following formulas evaluate all coefficients of that
same $b$.

\subsection{An explicit finite coefficient formula with the full arithmetic unit}

For each original root define the actual Taylor coefficients
\[
 w_{z,n}=[x^n]\frac{v_h(z+x)}{h_z(z+x)},\quad 0\leq n<m,
 \qquad
 \Psi_z(t)=\sum_{a=0}^{m-1}\frac{w_{z,m-1-a}}{a!}t^a.
 \tag{FTR.4}
\]
All these coefficients are specified in the original coordinates.
For an entirely finite evaluation, put
$v_{z,n}=v_h^{(n)}(z)/n!$ and
$h_{z,n}=h_z^{(n)}(z)/n!$. The product identity gives
\[
 w_{z,0}=\frac{v_{z,0}}{h_{z,0}},\qquad
 w_{z,n}=\frac{v_{z,n}-\sum_{d=1}^n h_{z,d}w_{z,n-d}}{h_{z,0}}
 \quad(1\leq n<m).
 \tag{FTR.5}
\]
The denominator is the nonzero original $h_z(z)$. This recurrence
retains every derivative, factorial, sign and complex phase of the
original unit; it introduces no additional source or normalization.

Define the polynomial in three formal variables
\[
 \mathcal P(X,Y,t)=\sum_{\epsilon,\eta=0}^1
                         X^\epsilon Y^\eta\Psi_{z_{\epsilon\eta}}(t),
 \qquad
 c_{ab,d}=d!\,[X^aY^bt^d]\mathcal P(X,Y,t)^k,
 \quad 0\leq d\leq r.
 \tag{FTR.6}
\]
Then the full local remainder of the comparison multiplier is exactly
\[
 \boxed{\quad
 b(\lambda_{ab}+x)
   =\chi_{ab}(\lambda_{ab}+x)
                  \sum_{d=0}^{r}c_{ab,d}x^{r-d}\pmod{x^e}.
 \quad}\tag{FTR.7}
\]
The right side, computed separately at all centers and assembled by
the unchanged Chinese remainder map, determines the unique original
class $b\in\mathcal E$. No coefficient system of unspecified rank
remains to be inverted.

Here is a proof of every coefficient in FTR.7. The full-unit formula
TWB42a says that, for every polynomial $P$,
\[
 B(1,[P])=(L_h^{\otimes k})
       \left(v_h^{\otimes k}P(s_1+\cdots+s_k)\right),
 \qquad L_h(f)=\sum_z\operatorname{Res}_{s=z}\frac{f(s)}{h(s)}ds.
 \tag{FTR.8}
\]
There is one full factor $v_h$ in each variable: the two factors
from the original arithmetic injection cancel only one factor with
the original inverse-unit residue, as already proved in TWB42a.
In an ordered root tuple $(z_1,\ldots,z_k)$, taking the residues in
the original order is exactly
\[
 [x_1^{m-1}\cdots x_k^{m-1}]
       P\left(\lambda+\sum_i x_i\right)
       \prod_{i=1}^k\frac{v_h(z_i+x_i)}{h_{z_i}(z_i+x_i)},
 \qquad \lambda=\sum_i z_i.
 \tag{FTR.9}
\]
Terms of total degree above $r$ in $\sum_i x_i$ cannot contribute.
Taylor expansion of $P$ therefore gives, at total degree $d$, the
coefficient
\[
 \frac{P^{(d)}(\lambda)}{d!}
 \sum_{\substack{a_1+\cdots+a_k=d\\0\leq a_i<m}}
        \frac{d!}{\prod_i a_i!}\prod_i w_{z_i,m-1-a_i}.
 \tag{FTR.10}
\]
Expansion of $\mathcal P^k$ enumerates each ordered root tuple once;
its $X,Y$ exponents are exactly the numbers of positive real and
imaginary root choices. Its $t^d$ coefficient enumerates the displayed
nilpotent exponents with the factors $1/a_i!$. Consequently the sum
of FTR.9 over every tuple is
\[
 B(1,[P])=\sum_{a,b=0}^k\sum_{d=0}^r
                   c_{ab,d}\frac{P^{(d)}(\lambda_{ab})}{d!}.
 \tag{FTR.11}
\]
On the other hand, the same value in FTR.3 is
\[
 L_\chi(bP)=\sum_{a,b}\sum_{d=0}^r
  \frac{P^{(d)}(\lambda_{ab})}{d!}
  [x^{r-d}]\frac{b(\lambda_{ab}+x)}{\chi_{ab}(\lambda_{ab}+x)}.
 \tag{FTR.12}
\]
All jets through order $r$ at the distinct centers can be prescribed
independently: the ideals $(S-\lambda_{ab})^e$ are pairwise comaximal,
and their Chinese remainder map is onto. Equality of FTR.11 and
FTR.12 for every $P$ therefore equates each displayed coefficient.
Multiplication by the retained local unit $\chi_{ab}$ proves FTR.7.
This also proves directly that the coefficient formula calculates
the same $b$ as TWB42, rather than a different pairing.

\subsection{The exact kernel and rank on every primary factor}

At a fixed center put
\[
 d_{ab}=\max\{d\in\{0,\ldots,r\}:c_{ab,d}\ne0\},
 \quad\hbox{with }d_{ab}=-1\hbox{ if every }c_{ab,d}=0.
 \tag{FTR.13}
\]
If $d_{ab}\geq0$, formula FTR.7 writes the local multiplier as
$x^{r-d_{ab}}$ times a unit: the leading coefficient of the last
factor is the nonzero number
$\chi_{ab}(\lambda_{ab})c_{ab,d_{ab}}$. If $d_{ab}=-1$, that
multiplier is zero. In the original local algebra
$\mathcal E_{ab}=\mathbb C[x]/x^e$ this gives the complete formulas
\[
 \begin{gathered}
 \ker(M_b|_{\mathcal E_{ab}})=x^{d_{ab}+1}\mathcal E_{ab},\qquad
 \operatorname{im}(M_b|_{\mathcal E_{ab}})
             =x^{r-d_{ab}}\mathcal E_{ab},\\
 \operatorname{rank}(M_b|_{\mathcal E_{ab}})=d_{ab}+1,
 \qquad
 \dim\ker(M_b|_{\mathcal E_{ab}})=r-d_{ab}.
 \end{gathered}\tag{FTR.14}
\]
For $d_{ab}=-1$, the image exponent is $e$, the kernel exponent
is zero, and these formulas give the entire kernel and zero image.
For every other value, multiplying the monomial basis by
$x^{r-d_{ab}}$ proves the formulas term by term; the remaining
unit is invertible and commutes with this multiplication.

Since $R_\chi$ is nondegenerate, FTR.3 identifies the right radical
of $B$ with $\ker M_b$. Both $B$ and $R_\chi$ are skew-Hermitian:
the original single-factor residue is skew-Hermitian and $k$ is odd,
while TWB40 proves the assertion for $R_\chi$. The left radical is
therefore the same subspace. In particular its exact total dimension
is $\sum_{a,b}(r-d_{ab})$, and its complementary rank is
$\sum_{a,b}(d_{ab}+1)$. These are finite explicit polynomial
coefficient calculations in FTR.4--6, with possible cancellations
retained rather than forbidden.

\subsection{Every perimeter block survives with its full nilpotent order}

Put the four nonzero original numbers
\[
 A_0=w_{z_{00},0},\qquad B_0=w_{z_{10},0},\qquad
 C=w_{z_{01},0},\qquad D_0=w_{z_{11},0},\qquad
 L(X,Y)=A_0+B_0X+CY+D_0XY.
 \tag{FTR.15}
\]
The subscripted $A_0,B_0,D_0$ in this display are scalar names, not
the boundary band or native metric used in the preceding sections.
The highest $t$ coefficient in each $\Psi_z$ is
$w_{z,0}/(m-1)!$. Hence the highest jet in FTR.6 is
\[
 c_{ab,r}=\frac{r!}{((m-1)!)^k}[X^aY^b]L(X,Y)^k.
 \tag{FTR.16}
\]
Every edge coefficient has a single noncancelling binomial expression:
\[
 \begin{array}{ll}
 [X^0Y^b]L^k=\binom{k}{b}A_0^{k-b}C^b,&
 [X^kY^b]L^k=\binom{k}{b}B_0^{k-b}D_0^b,\\[3pt]
 [X^aY^0]L^k=\binom{k}{a}A_0^{k-a}B_0^a,&
 [X^aY^k]L^k=\binom{k}{a}C^{k-a}D_0^a.
 \end{array}\tag{FTR.17}
\]
For example, the first coefficient can only choose factors with no
$X$, so is the stated coefficient of $(A_0+CY)^k$. The other three
follow by selecting all $X$, no $Y$, or all $Y$, respectively.
Every factor is nonzero and the original field has characteristic
zero. Thus all coefficients in FTR.17, including their unchanged
factor in FTR.16, are nonzero.

There are exactly $4(k+1)-4=4k$ distinct perimeter centers with
$a\in\{0,k\}$ or $b\in\{0,k\}$. At each of them $d_{ab}=r$,
so $b$ is a unit and its local comparison retains all $e$ jets.
The complete formula for an interior coefficient, where cancellation
must still be retained, is
\[
 [X^aY^b]L^k
  =\sum_{n=\max(0,a+b-k)}^{\min(a,b)}
    \frac{k!\,A_0^{k-a-b+n}B_0^{a-n}C^{b-n}D_0^n}
    {(k-a-b+n)!(a-n)!(b-n)!n!}.
 \tag{FTR.18}
\]
Indeed $n$ counts choices of $D_0XY$; the other three counts
are forced by the $X,Y$ exponents and total $k$. This proves the
formula without selecting phases or discarding cross terms. When
an interior highest coefficient vanishes, FTR.6 and FTR.13--14
calculate exactly which lower jets remain.

Let $\mathcal E_{\partial}$ be the direct sum of the $4k$ perimeter
primary factors under the original idempotents. It is invariant
under $A$ and under the original reflection: that reflection sends
$(a,b)$ to $(k-a,b)$. Residue duality pairs a primary factor only
with this reflected factor, since products of the other local
idempotents vanish. Thus $R_\chi$ restricts nondegenerately to
$\mathcal E_{\partial}$. The proved invertibility of $b$ on every
one of these factors now proves
\[
 \operatorname{rank}(B|_{\mathcal E_{\partial}})=4ke,
 \qquad
 \operatorname{inertia}(\mathrm iB|_{\mathcal E_{\partial}})
                         =(2ke,2ke).
 \tag{FTR.19}
\]
For the inertia assertion, no reflected factor is fixed because
$k$ is odd. In a pair of reflected factors the Hermitian matrix
of $\mathrm iB$ has the form
$\left(\begin{smallmatrix}0&C_0\\C_0^*&0\end{smallmatrix}\right)$,
with $C_0$ invertible by nondegeneracy. Congruence by
$\operatorname{diag}(I,C_0^{-1})$ changes this matrix to
$\left(\begin{smallmatrix}0&I\\I&0\end{smallmatrix}\right)$.
Its two subspaces $(u,u)$ and $(u,-u)$ have respectively positive
and negative forms, each of dimension $e$. There are $2k$ pairs,
giving exactly the inertia in FTR.19. Every center on this perimeter
has $2\Re\lambda_{ab}-k=2(2a-k)\delta\ne0$.

This explicitly locates a nondegenerate full-order part of the
original tensor-to-cyclic residue comparison even when its interior
is singular. The map is the original full-unit arithmetic injection
followed by the original tensor pairing; its recovered multiplier,
radical and perimeter are now fully calculated. Its nondegenerate
perimeter form is indefinite with the inertia just proved. The
positive boundary forms retain the nonzero action defect in TWB43--44;
neither this exact comparison nor its nondegeneracy sets that defect
to zero or establishes an RH endpoint.
